\begin{figure}[p]
\centering
\includegraphics[width=\textwidth]{supplementary/curated/anatomy_capacity_controls.pdf}
\caption{\textbf{Cohort, label and arbor-size controls bound modeled anatomical capacity.} \textbf{A}, Route-generated field capture versus budget in 47 disjoint same-mouse cells; morphology routes use 8.1\% of dense-feedback wiring at eight channels. Twenty modeled streams are averaged within each cell; bands/bars are 95\% cell-bootstrap intervals. \textbf{B}, Initial eight-cell cohort: capture restricted to directly typed presynaptic contacts. \textbf{C,D}, Second-mouse Pinky cohort: relative residual norm versus route budget, and candidate-site count versus effective field rank. Eligibility changes from ten cells at $K\leq4$ to eight at $K=8$; small arbors approach the rank ceiling. Panel C uses mean $\pm$ s.e.m., and its residual norm is not squared-energy capture. All fields are model generated. Main Fig.~7 equalizes a common broadcast and leads with actual-tree surrogate controls; these route-generated comparisons are structural positive controls, not measured biological teaching signals.}
\label{fig:si_anatomy_capacity_controls}
\label{fig:supp_v661}
\label{fig:supp_morphology_detail}
\label{fig:supp_pinky}
\end{figure}

\begin{figure}[p]
\centering
\includegraphics[width=\textwidth]{supplementary/curated/inhibitory_spatial_controls.pdf}
\caption{\textbf{Spatial matching limits apparent inhibitory targeting of ancestry domains.} \textbf{A}, Six layer-diverse reconstructions with mapped inhibitory contacts and physical scale bars. \textbf{B--D}, Observed-minus-matched tree distance, shared-path fraction and descendant-domain overlap in the twenty-target analysis, after aggregation within 92 presynaptic axons, and after joint path-distance/Euclidean matching. Symbols and intervals are means and 95\% bootstrap intervals over the displayed inferential unit. Only tree distance persists after axon aggregation; no endpoint persists after joint 3D matching. \textbf{E}, Descendant-input fraction by distal-dendrite- versus perisomatic-targeting class, retaining one contact per axonal clump. \textbf{F}, Inhibitory-contact domain size versus matched locations; no reliable overall localization advantage is detected. Cells derive from one mouse. These structural comparisons do not show that contacts carried teaching signals.}
\label{fig:si_inhibitory_spatial_controls}
\label{fig:supp_inhibitory}
\end{figure}

\begin{figure}[p]
\centering
\includegraphics[width=\textwidth]{supplementary/curated/coarse_energy.pdf}
\caption{\textbf{Irregular reconstructed trees concentrate modeled route-field energy coarsely, but the anatomy-specific excess does not replicate.} \textbf{A}, Weighted unbalanced tree-Haar contrasts recursively partition each reconstructed arbor into coarse, intermediate and fine supports, defined by descendant excitatory-weight fractions \(>1/4\), \(1/16\)--\(1/4\) and \(\leq1/16\), respectively. \textbf{B}, Non-scalar conductance-weighted ancestry-route energy fraction in the disjoint 47-cell cohort, compared with column-permuted ancestry and dimension-matched isotropic noise. Actual-route points and intervals are cell means and 95\% cell-bootstrap intervals; null markers are cell means. \textbf{C}, Signal-to-isotropic-noise power by scale for all 47 cells; diamonds show means and 95\% cell-bootstrap intervals. \textbf{D}, Actual-minus-permuted coarse energy in the original and disjoint cohorts. The positive original-cohort contrast is not reliable in the 47-cell cohort. Cells are the inferential units; all reconstructions derive from one mouse. Modeled route fields are not measured teaching signals.}
\label{fig:si_coarse_energy}
\label{fig:supp_irregular_wavelets}
\end{figure}

\begin{figure}[p]
\centering
\includegraphics[width=\textwidth]{supplementary/curated/anatomy_preprocessing.pdf}
\caption{\textbf{Label coverage and geometric preprocessing affect anatomical route dictionaries.} All analyses use eight initial reconstructed cells from one mouse. \textbf{A}, Direct-presynaptic versus target-proxy excitatory/inhibitory (E/I) labels for 2,012 jointly labeled mapped contacts; counts and row percentages expose the asymmetric disagreement. Contacts are nested observations. \textbf{B}, Direct-label fraction by compartment: individual cells and cell means with 95\% cell-bootstrap intervals. Label availability is not assumed random. \textbf{C}, Jaccard overlap of the four selected routes with the nominal hybrid-label, $5\,\mu$m dictionary across mapping thresholds and label choices. \textbf{D}, Capture of the fixed nominal model-generated field under heterogeneous log-radius perturbations, indexed by their standard deviation (SD); five draws per nonzero dose are averaged within cell. Curves and bars in \textbf{C,D} show means and 95\% cell-bootstrap intervals (5,000 draws). \textbf{E}, Distribution of the mean-radius compressed axial resistance divided by the exact series sum over raw edges. Length preservation does not imply axial-resistance preservation. \textbf{F}, Paired cell capture under mean-radius and exact series-resistance normalizations, with leak terms held fixed. The diagonal marks equality. Across 1,056 retained sensitivity conditions, target fields, probe sites and weights remain those of the nominal model. These are preprocessing sensitivities, not calibrated geometric uncertainties or tests of endogenous task alignment.}
\label{fig:si_anatomy_preprocessing}
\label{fig:supp_morphology_uncertainty}
\end{figure}
